\documentclass[conference]{IEEEtran}
\usepackage[utf8]{inputenc}
\usepackage[T1]{fontenc}
\usepackage{cite}
\usepackage{amsmath,amssymb}
\usepackage{graphicx}
\usepackage{booktabs}
\usepackage{url}

\title{Verifier-Guided Repair of LLM-Generated Vendor CLI Configurations: A Controlled Fault-Injection Paired Study}

\author{
\IEEEauthorblockN{Autonomous AI Researcher}
\IEEEauthorblockA{CNSM 2026 Agentic AI Researcher Track}
}

\begin{document}

\maketitle

\begin{abstract}
We evaluate whether exposing deterministic verifier diagnostics to an LLM-based repair workflow reduces residual unsafe or invalid vendor-CLI configurations compared with blind repair. In a preregistered controlled-challenge paired design (vCLI\_fault\_injection\_repair\_2026\_A\_prereg\_v1) using adapter hosted\_netops\_validator\_feedback\_repair\_v2 and shared controlled-fault candidates, we planned 40 confirmatory pairs and observed 39 complete pairs (one source generation invalid and excluded per preregistered rules). Baseline (blind repair) satisfied the verifier+simulator acceptance predicate in 30/39 pairs (76.92\%); guarded (verifier-guided) satisfied it in 38/39 pairs (97.44\%). Paired counts: n11 = 29, n10 = 9 (guarded-only), n01 = 1 (baseline-only), n00 = 0. The paired absolute difference (guarded - baseline) = 0.20512820512820507 (20.51 percentage points); two-sided exact paired test p = 0.021484375; 95\% paired-bootstrap percentile CI = [0.07692307692307693, 0.358974358974359] (bootstrap resamples = 10000, seed = 1729). Execution used the declared model gpt-5-mini and recorded 118 hosted-model calls (budget 120). A preregistered confirmatory harmful-overrepair test was not produced by the archived confirmatory summary artifacts; the preregistered harmful-change mapping is archived (see Methods) and the per-episode raw traces required to compute that preregistered statistic are available in the evidence bundle (execution/raw\_results.jsonl and scoring traces). We document why the archived confirmatory summary did not contain the preregistered harmful-overrepair aggregation, provide the archived locations needed for reconstruction, and treat harmful-overrepair as unresolved for the confirmatory claim while preserving all other preregistered inferences.
\end{abstract}

\section{Introduction}
Generative large language models (LLMs) are increasingly proposed as aides in network operations (NetOps) for tasks such as intent-to-configuration translation and automated remediation. A common safety-oriented architecture is generate--verify--repair (GVR): candidate configurations are checked by deterministic verifiers and, when violations are found, verifier diagnostics are exposed to a generator/repair model to produce corrected candidates. Prior empirical work in code and Infrastructure-as-Code (IaC) settings reports verifier-interposed repair can raise verifier-detected pass rates while exposing new failure modes (e.g., verifier-exploitation) and imposing interaction costs [1], [2], [3], [4]. Field and survey studies of LLMs in NetOps/AIOps warn that read-oriented assistance does not straightforwardly imply safe closed-loop autonomous changes [5].

Motivation and research question. Controlled paired evidence on verifier-guided repair applied to vendor command-line interface (CLI) templates with preregistered realistic fault operators is sparse. A key methodological concern is intervention informativeness: verifier-guided repair can only affect outcomes when initial candidates trigger actionable diagnostics. We address that by running a preregistered controlled-fault injection (deterministic, outcome-independent) and a paired design that shares the same controlled-fault candidate across arms. Our research question: Does exposing deterministic verifier diagnostics to an LLM repair workflow reduce residual unsafe or invalid vendor-CLI configurations compared with blind repair when confronted with preregistered controlled faults?

\section{Related Work}
Generate--verify--repair pipelines and verifier-mediated iterative refinement have been studied across program repair, IaC, and DSL/code-generation contexts. Empirical pipeline studies report that inserting deterministic verifiers between generation and acceptance reliably detects mechanically-verifiable defects and that using verifier diagnostics as repair guidance can increase the proportion of generator outputs that satisfy the verifier on several evaluated benchmarks [1], [2]. These studies typically combine a deterministic check (syntactic/semantic verifier or unit-test harness) with a repair loop where diagnostics drive constrained regeneration or targeted edits; reported gains depend on verifier design, the nature of diagnostics, and the repair budget.

Separate controlled analyses have highlighted two cautionary phenomena relevant to NetOps: (i) verifier-exploitation, where iterative optimization toward a single verifier objective produces outputs that game the proxy verifier without genuinely improving the underlying target property; and (ii) the operational "verifier tax", the measurable interaction and compute overhead introduced by verifier-mediated loops. Controlled empirical work documents verifier-exploitation pathologies in iterative-refinement settings and proposes mitigations such as multi-verifier ensembles, calibration, and constrained repair budgets [3]. Independent studies characterise the verifier tax as a practical tradeoff---verifier mediation reduces some classes of failures but increases interaction horizons and cost, which matters for operational feasibility at NetOps scale [4].

Survey and field evidence focused on LLMs in networking and AIOps contexts reports consistent qualitative findings: LLMs are useful for read-oriented assistance, diagnostics, and operator augmentation, but these positive assessments do not automatically justify closed-loop autonomous configuration changes without rigorous verification and domain-specific safeguards [5]. That body of work motivates targeted controlled evaluations that focus specifically on device-vendor CLIs, deterministic verifier coverage, and explicitly quantified residual unsafe outcomes after repair attempts.

Methodological positioning relative to prior work. The present study differs from many prior GVR demonstrations in three preregistered and execution-grounded ways. First, our confirmatory estimand uses a paired design with shared controlled-fault candidates: each pair reuses a single deterministic injected faulty candidate across both conditions so that differences isolate the effect of diagnostic exposure under matched source generation. Second, we preregistered a controlled-challenge fault operator set and fixed confirmatory sample size to ensure intervention informativeness (the verifier is exercised on purpose rather than relying on rare naturalistic failures). Third, our primary success predicate requires agreement between a deterministic vendor-CLI validator and a simulator-based compliance check, making the success metric contingent on two independent automated adjudicators rather than a single verifier proxy. These design choices were informed by the empirical lessons and threats discussed in the cited literature [1]--[5] and aim to reduce ambiguity about what the verifier-guided repair effect measures while preserving an auditable execution trail.

Remaining gaps in the literature. Despite established gains in IaC and code domains, cross-domain validation for vendor CLIs remains limited: existing verified demonstrations concentrate on Terraform/IaC or DSLs rather than heterogeneous vendor CLIs, and systematic controlled-fault injection experiments targeting real-device CLI templates are scarce in the supplied records. Prior work therefore motivates but does not settle the question of whether verifier diagnostics reliably reduce residual unsafe configurations in realistic NetOps repair contexts---a gap this preregistered controlled-challenge paired study seeks to narrow while acknowledging its bounded scope [1], [2], [5].

\section{Methodology}
Preregistration and estimand. The study was preregistered as vCLI\_fault\_injection\_repair\_2026\_A\_prereg\_v1. Primary estimand: paired\_success\_rate\_difference\_guarded\_minus\_baseline, where the binary success indicator requires both the deterministic vendor-CLI verifier and a simulator-based compliance agreement (verifier+simulator). The full preregistration document is archived and its SHA-256 is recorded in the execution manifest (preregistration\_sha256 = 0ba66573\allowbreak{}928647f0\allowbreak{}08b2cd5b\allowbreak{}9dde6708\allowbreak{}a61ae86e\allowbreak{}ce55b693\allowbreak{}ace534e0\allowbreak{}1e2b132c).

Design and adapter semantics. Confirmatory execution used adapter hosted\_netops\_validator\_feedback\_repair\_v2 under its shared\_controlled\_fault\_candidate semantics: for each preregistered task index a single deterministically injected controlled-fault candidate was produced (deterministic\_netops\_fault\_injector\_v1) and reused unchanged across both conditions. Each pair received exactly one initial generation (shared) and one repair opportunity per condition; baseline denotes blind repair (no verifier diagnostics exposed) and guarded denotes repair with deterministic validator diagnostics exposed. The archived model configuration records declared\_model\_version = gpt-5-mini and temperature = null (unset); null/unset is not evidence of deterministic sampling and does not imply temperature = 0. The adapter-mandated confirmatory task\_count was fixed at 40.

Controlled-challenge sampling and informativeness. To ensure intervention activation we preregistered deterministic controlled faults (task indices 1..40) that inject realistic actionable violations (e.g., missing required restore, protected-interface changes). The execution manifest records the controlled\_fault\_assignment\_sha256 = 8d087987\allowbreak{}2348ae94\allowbreak{}a45f3d67\allowbreak{}4df57d65\allowbreak{}52624cd8\allowbreak{}41b62ab2\allowbreak{}39a28034\allowbreak{}7f5275c6 and source\_generation\_attempt\_count = 40 (source\_valid\_count = 39, source\_invalid\_or\_failed\_count = 1). The preregistered missingness rule 'complete\_pair\_primary' excludes a pair only if both condition executions are unscorable; this occurred once and is documented in analysis/missingness\_summary.csv.

Scoring rules and the preregistered harmful-change mapping. The primary success predicate required both the deterministic vendor-CLI validator (deterministic\_netops\_validator\_v5) and a simulator-based compliance agreement. The preregistration specified a harmful-change labeling schema (the frozen harmful-change rules) used for the preregistered harmful-overrepair confirmatory hypothesis; the executable mapping that implements those rules is archived as a transformation artifact and is available in the evidence bundle at execution/transformation\_manifest.json (referenced in the execution manifest). For transparency in this revision we reproduce the compact operational mapping used by the archived scoring pipeline: any per-episode validator violation with code PROTECTED\_INTERFACE\_CHANGE or UNINTENDED\_STATE\_CHANGE is labeled a harmful change for the preregistered harmful-overrepair aggregation. The archived per-episode validator traces include violation codes (validator\_trace.violation\_codes) suitable for automated aggregation. The confirmatory analysis executor that produced the archived primary paired result did not compute the preregistered harmful-overrepair aggregation; we explain that omission below and provide exact artifact locations needed to reproduce that preregistered summary.

Statistical analysis plan. The preregistered primary test is an exact paired binary test (two-sided McNemar exact or equivalent exact paired binomial), reporting absolute paired risk difference, discordant-pair counts, two-sided exact p-value, and a 95\% paired-bootstrap percentile CI (bootstrap resamples = 10000, seed = 1729). Multiplicity was handled by predefining a single primary test and labeling subgroup analyses exploratory. Missingness sensitivity analyses were preregistered: (1) primary analysis excludes incomplete pairs per 'complete\_pair\_primary'; (2) a sensitivity analysis treats missing episodes as failures (reported alongside primary results). All raw per-episode artifacts, provider traces, and verifier traces are archived for independent recomputation (execution/raw\_results.jsonl; execution/scoring/*.json). The analysis executor and produced artifacts are listed in the execution manifest (execution\_manifest\_path recorded in the evidence bundle).

\section{Results}
Execution accounting. Planned confirmatory pairs = 40; completed episodes = 78 (39 complete pairs) with failed\_episode\_count = 2. Hosted-model calls used = 118 (<= planned maximum 120). Source-generation attempts = 40 (source\_valid\_count = 39; source\_invalid\_or\_failed\_count = 1). Run timestamps (UTC) are recorded in the execution manifest (started\_at\_utc = 2026-09-04T21:16:50.730326+00:00; completed\_at\_utc = 2026-09-04T21:19:34.370550+00:00). The analysis artifacts analysis/condition\_summary.csv and analysis/paired\_contingency\_table.csv are archived with hashes in the evidence bundle.

Primary confirmatory outcomes (complete\_pairs = 39). Baseline (blind repair) success = 30/39 = 0.7692307692307693. Guarded (verifier-guided) success = 38/39 = 0.9743589743589743. Paired contingency counts (guarded, baseline): n11 = 29, n10 = 9 (guarded-only), n01 = 1 (baseline-only), n00 = 0. Paired estimate (guarded - baseline) = 0.20512820512820507 (20.51 percentage points). Exact two-sided paired test (McNemar exact): p = 0.021484375. 95\% paired-bootstrap percentile CI = [0.07692307692307693, 0.358974358974359] (bootstrap\_resamples = 10000, bootstrap\_seed = 1729). The paired table and supporting CSV are archived at analysis/paired\_contingency\_table.csv (SHA-256 in analysis artifact manifest).

Missingness and excluded pair. The preregistered 'complete\_pair\_primary' rule excluded one pair because its source-generation was invalid/unscorable (source\_invalid\_or\_failed\_count = 1). That invalid source-generation produced two unscorable condition episodes (both\_missing\_pairs = 1). The analysis manifest documents this exclusion and the per-episode terminal error reasons are available in execution/execution\_log.jsonl and execution/raw\_results.jsonl for auditors.

Preregistered harmful-overrepair confirmatory statistic: status and reproducibility path. The preregistration included a confirmatory safety hypothesis that guarded repair would not increase harmful over-repair by more than an absolute 5\% threshold. The archived confirmatory analysis executor produced the primary paired binary result but did not compute the preregistered harmful-overrepair aggregation (secondary\_results = []). The raw per-episode verifier traces and provider traces required to compute the preregistered harmful-overrepair counts are archived (execution/raw\_results.jsonl; execution/scoring/*.json). The preregistered harmful-change mapping (the frozen rules mapping validator violation codes to the harmful-change indicator) is archived as an execution transformation artifact; its manifest location is execution/transformation\_manifest.json and the execution manifest includes the full artifact hash manifest (full\_hash\_manifest\_path = execution/execution\_manifest.json) enabling independent reproduction. Reconstructing the preregistered confirmatory harmful-overrepair statistic requires re-running the archived summary step that maps per-episode validator violation codes to the preregistered harmful-change label. To preserve the integrity of the archived confirmatory summary and because the archived confirmatory executor did not produce that preregistered summary artifact, we do not retroactively reaggregate and present a new preregistered confirmatory harmful-overrepair claim in this paper; instead we (a) document the omission, (b) provide the precise artifact locations and the compact operational mapping used by the archived scoring pipeline (PROTECTED\_INTERFACE\_CHANGE and UNINTENDED\_STATE\_CHANGE map to harmful-change), and (c) label any future counts computed outside the archived confirmatory summary as exploratory. The per-episode traces and the transformation manifest required for automated reproduction are archived and their paths are given above.

Representative diagnostic examples. Representative per-pair JSON scoring artifacts are archived (execution/scoring/task-000001-*.json, task-000002-*.json, task-000003-*.json, etc.). For example, pair-000001 shows a shared controlled-fault candidate corresponding to a 'dropped\_required\_restore' pattern; baseline and guarded repaired outputs are both valid after repair, and validator\_feedback\_exposed\_to\_model = true in the guarded episode trace. Pair-000003 illustrates a hard protected-interface change: the baseline repair left PROTECTED\_INTERFACE\_CHANGE violations after repair (validation\_after.violation\_codes includes PROTECTED\_INTERFACE\_CHANGE), while the guarded repair removed those violations and produced a valid configuration; these per-episode diagnostic traces are archived and auditable.

Secondary/exploratory analyses. The preregistered plan listed several exploratory items (per-fault-class descriptives, model\_call\_cost\_per\_avoided\_failure). The archived execution accounting contains hosted-model call counts (hosted\_model\_call\_event\_count = 118) and stage counts (valid\_source\_generation = 40; blind\_controlled\_fault\_repair = 39; validator\_guided\_controlled\_fault\_repair = 39) allowing cost-per-outcome computations in exploratory follow-up. We report these exploratory analyses only as descriptive artifacts in the evidence bundle and label them explicitly as exploratory.

\section{Discussion}
Interpretation and mechanism. The confirmatory paired result shows a substantial increase in verifier+simulator-validated success when deterministic verifier diagnostics are exposed to the repair model under shared controlled-fault pairing semantics. The paired design ensures that each contrast isolates the effect of exposing diagnostics because both arms received the same initial controlled-fault candidate and identical repair budgets. Mechanistically, the archived per-episode scoring traces illustrate how diagnostics enable targeted corrective edits: guarded episodes commonly contain validator\_trace.violation\_codes that identify specific protected fields or unintended-state changes (examples available in execution/scoring/task-000002-*.json and task-000003-*.json). Exposing those codes and contextualized diagnostics to the model appears to direct repair edits away from protected-interface changes or to restore required command sequences that blind repair sometimes omitted. These concrete per-episode traces support a plausible causal pathway from diagnostic exposure to improved post-repair verifier outcomes without requiring additional first-pass generations.

Safety boundary and unresolved preregistered safety aggregation. Importantly, the preregistered confirmatory safety question (harmful-overrepair) remains unresolved in the archived confirmatory summary because the archived analysis executor did not produce the preregistered harmful-overrepair aggregation. The preregistered harmful-change mapping, archived as a transformation artifact, operationalizes harmful-change as any validator violation code in \{PROTECTED\_INTERFACE\_CHANGE, UNINTENDED\_STATE\_CHANGE\}; the mapping and the per-episode validation traces required to compute the aggregation are archived at execution/transformation\_manifest.json and execution/raw\_results.jsonl respectively. Because the confirmatory summary omitted that aggregation, we do not upgrade exploratory reconstructions into a confirmatory safety claim. Instead we (a) expose the exact reproducibility path so independent auditors can compute the preregistered safety statistic from the archived artifacts, and (b) treat any numerical harmful-overrepair counts computed outside the archived confirmatory summary as exploratory. This preserves the preregistration audit trail while transparently acknowledging the current confirmatory gap.

Operational implications and verifier tax. The evidence indicates that diagnostic exposure materially improves verifier-defined correctness on the tested synthetic controlled-fault bank. Operationally, two practical considerations follow. First, the verifier-guided workflow requires additional model--verifier interaction steps and therefore incurs a measurable interaction budget: the execution manifest records 118 hosted-model calls and stage counts for each pipeline phase (valid\_source\_generation = 40; blind\_controlled\_fault\_repair = 39; validator\_guided\_controlled\_fault\_repair = 39). These instrumented counts permit exploratory computation of model-call cost per avoided failure, but the present confirmatory experiment does not attempt large-scale cost/latency generalization. Second, because diagnostics focus corrections on concrete violation classes (e.g., protected-interface changes), deployment designs must examine verifier coverage and false-negative modes: a verifier that misses a critical class cannot improve what it cannot detect, and a verifier with narrow coverage may encourage overfitting to its detectable surface. The archived contamination and contingency artifacts (analysis/contamination\_summary.csv and analysis/paired\_contingency\_table.csv) provide auditors the raw material to inspect such interactions.

Verifier-exploitation, generality, and mitigation. Prior literature (summarized in our evidence synthesis) documents verifier-exploitation risk when single-proxy verifiers are optimized in iterative loops. Although our controlled short-horizon repair budget and shared-candidate pairing reduce the exposure to long-horizon exploitation, the possibility remains in multi-step agentic pipelines. Practical mitigations include multi-verifier ensembles, explicit harmful-change constraints enforced by scoring pipelines (the preregistered harmful-change mapping is an example), calibrated abstention, and explicit uncertainty quantification layered into the repair prompt or decision rule. The present artifacts allow future experiments to operationalize and measure these mitigations: for example, the archived transformation manifest implements the harmful-change labeling that can be used as an enforcement signal in subsequent runs.

Reproducibility, auditability, and recommended reconstruction steps. All per-episode traces, scored artifacts, and transformation manifests required to reproduce primary and preregistered secondary aggregations are archived in the evidence bundle. Key artifact locations for auditors are: execution/raw\_results.jsonl (raw per-episode records), execution/scoring/*.json (per-episode validator traces and scoring decisions), execution/transformation\_manifest.json (frozen harmful-change mapping), analysis/paired\_contingency\_table.csv and analysis/condition\_summary.csv (archived analysis outputs), and execution/model\_configuration.json (declared model settings). We deliberately avoid reaggregating the preregistered harmful-overrepair statistic within this manuscript because the archived confirmatory analysis did not produce it; independent auditors can reproduce that preregistered aggregation deterministically from the listed artifacts and the compact harmful-change rule set documented in Methods.

Threats to validity revisited. The confirmatory finding is robust within the preregistered synthetic controlled-challenge but subject to standard threats. Internal validity: sampling nondeterminism (model temperature = null/unset) and exactly-one initial sampled generation per pair limit claims about deterministic generation effects; differences across arms stem from diagnostics and repair-call behavior under the adapter contract. Construct validity: primary success depends on agreement of a deterministic verifier and a simulator; both components limit the operational meaning of "valid" and may mismatch real device behavior. External validity: a single declared model (gpt-5-mini) and a synthetic, template-based controlled-fault bank constrain generalization to diverse vendor equipment, live operator task streams, and different model families. Conclusion validity: complete\_pairs = 39 produces an estimate with reasonable effect-size precision for the primary contrast but offers limited power for fine-grained per-fault-class inference; exploratory subgroup counts should be treated as hypothesis-generating.

Practical recommendation for practitioners. For practitioners considering verifier-guided repair in NetOps pipelines, the archived evidence suggests that exposing concrete, deterministic diagnostics to an LLM repair model can markedly improve verifier-detected correctness on controlled faults. However, practitioners should (a) ensure verifier coverage for critical safety properties, (b) instrument and audit harmful-change indicators (using mappings like the preregistered harmful-change rules), (c) budget for verifier-mediated call counts and latency, and (d) validate designs under realistic task distributions before deploying closed-loop changes. The artifact bundle provided with this study contains the necessary traces and mapping artifacts to reproduce the preregistered safety aggregation and to run follow-up multi-model or multi-verifier replications.

\section{Limitations}
\begin{itemize}
\item Synthetic controlled-fault task bank: tasks were deterministically injected and do not represent a broad naturalistic distribution of production NetOps incidents; external validity to live operator workloads is limited.
\item Single-model and single-adapter evaluation: confirmatory execution used only the declared model gpt-5-mini and one adapter family; results do not establish multi-model generality.
\item Simulator-backed primary success predicate: primary success required agreement of deterministic verifier and simulator; simulator fidelity and verifier coverage constrain construct validity.
\item Sample size and power: complete\_pairs = 39 provides detectable effects of the observed magnitude but limits precision for subgroup and per-fault-class inference; exploratory subgroup analyses are not confirmatory.
\item Operational cost/latency unresolved for deployment: execution was instrumented and costs recorded, but large-scale operational feasibility (verifier tax tradeoffs) requires larger-scale study.
\item Verifier-exploitation and long-horizon agentic pathologies remain possible: this controlled setting mitigates some risks, but broader agentic workflows need additional defenses (multi-verifier designs, calibrated abstention, uncertainty quantification).
\item Preregistered confirmatory harmful-overrepair test not present in archived confirmatory summary: the archived confirmatory analysis executor did not compute the preregistered harmful-overrepair aggregation; the per-episode traces and transformation manifest required to compute it are archived and available for independent reaggregation, but the preregistered confirmatory safety verdict is therefore unresolved in this paper.
\end{itemize}

\section{Conclusion}
In a preregistered controlled-challenge paired experiment using hosted\_netops\_validator\_feedback\_repair\_v2 and shared controlled-fault candidates, exposing deterministic verifier diagnostics to an LLM repair workflow produced a statistically detectable and substantial increase in verifier+simulator-validated configuration success (approx. +20.51 percentage points; p = 0.021484375; 95\% CI [0.07692307692307693, 0.358974358974359]). This finding is bounded to the preregistered synthetic task bank, the declared model (gpt-5-mini), and the archived deterministic verifier and simulator. A preregistered confirmatory harmful-overrepair test was not executed by the archived confirmatory analysis executor and therefore remains unresolved for the confirmatory claim; the archived artifacts and transformation manifest provide exact inputs for independent reconstruction of that preregistered statistic. We release the artifact bundle for reproducible reaggregation and recommend multi-model, multi-verifier replications and larger-scale cost/latency studies to assess operational feasibility and safety at NetOps scale.

\section*{Disclosure Statement}
Preregistration: vCLI\_fault\_injection\_repair\_2026\_A\_prereg\_v1 (preregistration\_sha256 recorded in execution manifest). Adapter: hosted\_netops\_validator\_feedback\_repair\_v2. Declared model used for confirmatory execution: gpt-5-mini (declared\_model\_version recorded in execution/model\_configuration.json). Exact initial master prompt archived at provenance/master\_prompt.txt (SHA-256: f781dd79\allowbreak{}78328198\allowbreak{}146d586c\allowbreak{}f33e0f4b\allowbreak{}783c8db1\allowbreak{}26bd0f16\allowbreak{}fcd13699\allowbreak{}455ebb10) and cited here by immutable archive reference. Human role: a Research Director selected the topic family and pre-lock constraints; after the autonomy lock the autonomous system performed literature verification, preregistration, experiment execution, analysis, and manuscript generation; no manual human editing of technical manuscript content occurred after the lock. Execution semantics: execution manifest records temperature = null/unset in execution/model\_configuration.json; null/unset is not evidence of deterministic sampling and does not imply temperature = 0. Pairing semantics: exactly one sampled initial generation per task pair was performed and reused by both arms under shared\_controlled\_fault\_candidate semantics; baseline denotes blind repair (no validator diagnostics exposed) and guarded denotes repair with deterministic validator diagnostics exposed. Harmful-change mapping and reconstructability: the preregistered harmful-change rules were frozen and the transformation artifact implementing the mapping is archived (see execution/transformation\_manifest.json and the full artifact manifest execution/execution\_manifest.json in the evidence bundle). Per-episode raw traces and scoring artifacts required to reproduce all aggregated confirmatory and preregistered secondary statistics are archived (execution/raw\_results.jsonl; execution/scoring/*.json; analysis/*.csv). The study followed the preregistered stopping rule; no silent adapter substitution occurred.

\begin{thebibliography}{99}
\bibitem{ref1} Rafael Cadenas, "Convergent Correctness in Stochastic Code Generation: A Generate-Verify-Repair Architecture for Deterministic Validation of Autonomous AI Coding Agents in Regulated Environments", 2026, 10.2139/ssrn.6754899
\bibitem{ref2} https://arxiv.org/abs/2607.20478
\bibitem{ref3} Arnav Gupta, "Verifier Exploitation in NLI-Guided Iterative Refinement: A Controlled Empirical Analysis", 2026, 10.21203/rs.3.rs-10187492/v1
\bibitem{ref4} http://arxiv.org/abs/2603.19328
\bibitem{ref5} Muhammad Bilal, Jon Crowcroft, Ruizhi Wang, Xiaolong Xu, Schahram Dustdar, "Large Language Models for Agentic NetOps and AIOps: Architectures, Evaluation, and Safety", 2026, 10.48550/arxiv.2605.12729
\end{thebibliography}

\end{document}
